\section{Phân bón và nguyên tố dinh dưỡng}

\subsection{Vai trò của phân bón}

Cây trồng cần nhiều nguyên tố dinh dưỡng để sinh trưởng. Trong đó, nitrogen, phosphorus và potassium là ba nguyên tố đa lượng thường được bổ sung qua phân bón hóa học.

\begin{center}
\begin{tabularx}{0.96\textwidth}{>{\bfseries}C{2.2cm}YC{3.5cm}}
\toprule
Nguyên tố & Vai trò chính & Biểu hiện thường gặp khi thiếu \\
\midrule
N & Thúc đẩy sinh trưởng thân, lá & Cây sinh trưởng kém, lá nhạt màu \\
P & Hỗ trợ phát triển bộ rễ, ra hoa, tạo hạt & Rễ phát triển kém, cây chậm lớn \\
K & Tăng khả năng chống chịu, góp phần nâng cao chất lượng nông sản & Cây yếu, khả năng chống chịu giảm \\
\bottomrule
\end{tabularx}
\end{center}

\begin{ghinho}{Các nhóm phân bón thường gặp}
\begin{itemize}
\item \textbf{Phân đạm}: cung cấp nitrogen; thường gặp \ce{NH4NO3}, \ce{NH4Cl}, \ce{(NH4)2SO4}, \ce{CO(NH2)2}.
\item \textbf{Phân lân}: cung cấp phosphorus; chứa các phosphate.
\item \textbf{Phân kali}: cung cấp potassium; thường gặp \ce{KCl}, \ce{K2SO4}.
\item \textbf{Phân NPK}: cung cấp đồng thời nitrogen, phosphorus và potassium.
\end{itemize}
\end{ghinho}

Khi chọn phân bón, không chỉ cần biết tên loại phân mà còn phải xem \textbf{thành phần dinh dưỡng} và \textbf{điều kiện đất, cây trồng}.

\subsection{Hàm lượng chất dinh dưỡng}

Cách tính hàm lượng một nguyên tố trong hợp chất không khác cách tính phần trăm khối lượng đã học:
\[
\%X=\frac{m_X\text{ trong }1\text{ mol chất}}{M_{\text{chất}}}\times100\%.
\]

\begin{vidu}
Tính phần trăm nitrogen theo khối lượng trong urea \ce{CO(NH2)2}. Cho $M_{\ce{CO(NH2)2}}=60$ g/mol.
\end{vidu}
\begin{loigiai}
Trong $1$ mol urea có $2$ mol nguyên tử nitrogen, nên:
\[
m_N=2\times14=\SI{28}{\gram}.
\]
Do đó:
\[
\%N=\frac{28}{60}\times100\%\approx46.7\%.
\]

\textbf{Vậy,} urea chứa khoảng $46.7\%$ nitrogen theo khối lượng.
\end{loigiai}

\begin{ghinho}{Hai chiều tính thường gặp}
\[
\boxed{m_{\text{phân bón}}}
\longrightarrow
\boxed{m_{\text{chất dinh dưỡng}}}
\]
\[
\boxed{m_{\text{chất dinh dưỡng cần}}}
\longrightarrow
\boxed{m_{\text{phân bón cần dùng}}}.
\]
Nếu cây chỉ sử dụng được một phần chất dinh dưỡng đã bón, phải tính thêm hiệu suất sử dụng.
\end{ghinho}

\begin{vidu}
Một ruộng cần được bổ sung thực tế \SI{46}{\kilo\gram} nitrogen. Dùng urea chứa $46\%$ N, giả sử cây sử dụng được $80\%$ lượng nitrogen đã bón. Tính khối lượng urea cần dùng.
\end{vidu}
\begin{loigiai}
Lượng nitrogen phải bón:
\[
m_N=\frac{46}{0.80}=\SI{57.5}{\kilo\gram}.
\]
Khối lượng urea cần dùng:
\[
m_{\text{urea}}=\frac{57.5}{0.46}=\SI{125}{\kilo\gram}.
\]

\textbf{Vậy,} cần dùng \SI{125}{\kilo\gram} urea.
\end{loigiai}

\subsection{Đọc nhãn phân NPK}

Trên nhiều loại phân NPK, ba số ghi trên bao bì cho biết hàm lượng quy ước của ba thành phần dinh dưỡng chính theo thứ tự N--P--K. Ví dụ, nhãn NPK 16--16--8 cho biết sản phẩm chứa $16\%$ N và các hàm lượng quy ước tương ứng của phosphorus và potassium theo cách ghi trên nhãn phân bón.

Khi đề bài cho trực tiếp phần trăm của từng thành phần, hãy sử dụng đúng số liệu của đề; không tự đổi sang phần trăm nguyên tố nếu đề không yêu cầu.

\begin{tuluyen}
Tính phần trăm nitrogen theo khối lượng trong các chất sau:
\begin{tasks}(2)
\task \ce{NH4NO3}.
\task \ce{NH4Cl}.
\task \ce{(NH4)2SO4}.
\task \ce{CO(NH2)2}.
\end{tasks}
\dapso{$35.0\%$; $26.2\%$; $21.2\%$; $46.7\%$.}
\begin{loigiaiBT}
Dùng công thức
\[
\%N=\frac{m_N\text{ trong }1\text{ mol chất}}{M_{\text{chất}}}\times100\%.
\]
Ta có:
\[
\%N_{\ce{NH4NO3}}=\frac{28}{80}\times100\%=35.0\%,
\]
\[
\%N_{\ce{NH4Cl}}=\frac{14}{53.5}\times100\%\approx26.2\%,
\]
\[
\%N_{\ce{(NH4)2SO4}}=\frac{28}{132}\times100\%\approx21.2\%,
\]
\[
\%N_{\ce{CO(NH2)2}}=\frac{28}{60}\times100\%\approx46.7\%.
\]

\textbf{Vậy,} trong bốn chất trên, urea có hàm lượng nitrogen lớn nhất.
\end{loigiaiBT}
\end{tuluyen}
